\chapter{L'Hopital's Rule and Indeterminate Forms}

Some limits answer themselves. If you want $\lim_{x\to 2}(x^2+1)$, you put in $2$ and
read off $5$; continuity does the work and there is nothing to think about. The
interesting limits are the ones that \emph{refuse} to answer themselves --- the ones
where substituting the target value produces not a number but a piece of nonsense
like $\tfrac{0}{0}$ or $\tfrac{\infty}{\infty}$. These are the \emph{indeterminate
forms}, and despite the menacing name they are the most important limits in all of
calculus. The derivative itself is one of them: $f'(a)=\lim_{h\to 0}\frac{f(a+h)-f(a)}{h}$
is a $\tfrac{0}{0}$ in disguise, every single time.

This chapter is about \emph{L'Hopital's rule}, which resolves
a huge class of these forms by a single, almost shocking idea: when a ratio collapses
to $\tfrac{0}{0}$, the limit is decided not by the values of the two functions (those
are both heading to zero) but by the \emph{rates} at which they get there. Replace the
functions by their derivatives and the fog lifts. We will see why that is true, where
it comes from, exactly when it is allowed, and the several ways people get it wrong.
And along the way we will return to a question from much earlier in the book --- which
functions grow faster than which --- and finally settle it cleanly.

\section{What ``indeterminate'' really means}

Start by getting the word right, because almost everyone misreads it. An
indeterminate form is \emph{not} an expression with no value. It is an expression
whose value is \emph{not determined by the form alone}. That is a precise and
different claim, and the whole subject hinges on it.

\begin{focusbox}
``Indeterminate'' is a statement about \emph{information}, not about
\emph{impossibility}. The symbol $\tfrac{0}{0}$ tells you that both the top and the
bottom are shrinking to zero --- but it does not tell you \emph{how fast}, and that is
the one thing that decides the limit. The form has thrown away the information you
need. Your job is to dig it back out.
\end{focusbox}

To feel why $\tfrac{0}{0}$ genuinely carries no answer, watch four limits that all
wear that same face as $x\to 0$, yet land in four different places:
\[
\frac{2x}{x}\to 2,\qquad
\frac{x}{x^2}\to \infty,\qquad
\frac{x^2}{x}\to 0,\qquad
\frac{x}{x}\to 1.
\]
In every one, plugging in $x=0$ gives $\tfrac{0}{0}$. If $\tfrac{0}{0}$ had a fixed
value, all four would have to share it. They do not: the answers are $2$, $\infty$,
$0$, and $1$. So $\tfrac{0}{0}$ cannot be a number. It is a \emph{question}, and the
answer depends entirely on the relative speed of the vanishing.

\begin{intuitionbox}
Think of two runners both decelerating to a dead stop at the finish line. The ratio
of their \emph{remaining distances} as they crawl in is $\tfrac{0}{0}$ --- both gaps
go to zero. But whether that ratio approaches $2$ or $\tfrac12$ or $1$ depends only on
their relative \emph{speeds} near the line. The runner who is still moving twice as
fast keeps twice the gap. The form $\tfrac{0}{0}$ hides the speeds; L'Hopital's rule
hands them back to you as derivatives.
\end{intuitionbox}

Contrast this with a form that is \emph{not} indeterminate. The expression
$\tfrac{1}{0}$ (a nonzero number over something shrinking to zero) is not
indeterminate: it blows up, with a sign you can usually pin down. And $\tfrac{0}{1}$
is not indeterminate either: it is just $0$. The trouble appears only when a
``smallness'' or a ``largeness'' competes against another of the same kind, so that
neither wins automatically.

\begin{focusbox}
\textbf{The seven indeterminate forms.} Up to the tricks of this chapter, all the
indeterminacy in elementary calculus comes in seven flavors:
\[
\frac{0}{0},\qquad \frac{\infty}{\infty},\qquad 0\cdot\infty,\qquad
\infty-\infty,\qquad 1^{\infty},\qquad 0^{0},\qquad \infty^{0}.
\]
The first two are ratios --- L'Hopital's rule attacks them directly. The middle two,
$0\cdot\infty$ and $\infty-\infty$, are \emph{disguised} ratios; we rewrite them as
$\tfrac{0}{0}$ or $\tfrac{\infty}{\infty}$. The last three are \emph{powers}; we take a
logarithm to drag them down into a product, then into a ratio. Everything reduces to
the first two.
\end{focusbox}

\begin{trapbox}
Do not put forms like $\tfrac{\infty}{0}$, $0^{\infty}$, $\infty+\infty$, or
$\tfrac{1}{0}$ on the indeterminate list. They are \emph{determinate}: $\tfrac{1}{0}$
diverges, $0^{\infty}=0$, $\infty+\infty=\infty$, $\tfrac{\infty}{0}$ diverges. The
test is always the same --- is there a genuine competition where the answer could come
out anywhere depending on the rates? If one side obviously wins, the form is not
indeterminate and you should just read off the answer.
\end{trapbox}

\section{The ratio-of-rates idea}

Here is the insight that makes the whole chapter. Slow down for it.
Suppose $f(a)=0$ and $g(a)=0$, and we want
\[
\lim_{x\to a}\frac{f(x)}{g(x)}.
\]
Near $x=a$ both functions are small. But ``small'' is not one number --- a smooth
function that passes through zero is, very close to that zero, almost exactly a
straight line: its tangent line. Concretely, for $x$ near $a$,
\[
f(x)\approx f(a)+f'(a)(x-a)=f'(a)\,(x-a),\qquad
g(x)\approx g(a)+g'(a)(x-a)=g'(a)\,(x-a),
\]
because $f(a)=g(a)=0$ kills the constant terms. Now form the ratio and watch the
common factor $(x-a)$ cancel:
\[
\frac{f(x)}{g(x)}\approx \frac{f'(a)\,(x-a)}{g'(a)\,(x-a)}=\frac{f'(a)}{g'(a)}.
\]
The vanishing $(x-a)$ in the numerator and denominator was the entire source of the
$\tfrac{0}{0}$ trouble --- and it cancels. What survives is the ratio of the two
\emph{slopes} at $a$. That is L'Hopital's rule, glimpsed through the linear
approximation.

\begin{happeningbox}
\textbf{What is really happening in $\tfrac{0}{0}$.} Both functions are diving toward
zero, and near their common zero each one is essentially its tangent line --- a number
(its slope) times the same small quantity $(x-a)$. The ratio of two such lines is the
ratio of their slopes; the small quantity $(x-a)$ cancels top and bottom and never
gets a chance to be ``zero over zero.'' So the limit is not really $\tfrac{0}{0}$ at
all. It is $\tfrac{f'(a)}{g'(a)}$ --- the ratio of the rates at which the two functions
leave zero. L'Hopital's rule is the act of looking past the values, which are both $0$
and tell you nothing, to the slopes, which decide everything.
\end{happeningbox}

Look back at the four warm-up limits with this lens. For $\tfrac{2x}{x}$ the slopes at
$0$ are $2$ and $1$, ratio $2$. For $\tfrac{x^2}{x}$ the numerator leaves zero with
slope $0$ (it is flat there) while the denominator has slope $1$, ratio $0$. For
$\tfrac{x}{x^2}$ it is $\tfrac10$, which is why that one blew up. The ratio of slopes
predicts each answer exactly. The picture below shows the mechanism: near the shared
zero, $f$ and $g$ are their tangent lines, and the quotient is the quotient of the
lines.

\begin{figure}[H]
\centering
\begin{tikzpicture}
\begin{axis}[
  width=11cm, height=6.5cm,
  axis lines=middle,
  xlabel={$x$}, ylabel={$y$},
  xtick={1}, xticklabels={$a$},
  ytick=\empty,
  xmin=-0.2, xmax=2.4, ymin=-1.2, ymax=1.8,
  samples=120,
  clip=false,
]
% f curve through (1,0) with slope ~0.7
\addplot[curveblue, very thick, domain=0.2:2.1] {0.7*(x-1)+0.5*(x-1)^2};
% g curve through (1,0) with slope ~0.35
\addplot[targetgreen, very thick, domain=0.2:2.1] {0.35*(x-1)-0.3*(x-1)^2};
% tangent lines at a
\addplot[curveblue, dashed, domain=0.45:1.7] {0.7*(x-1)};
\addplot[targetgreen, dashed, domain=0.45:1.7] {0.35*(x-1)};
\addplot[only marks, mark=*, mark size=1.6pt, black] coordinates {(1,0)};
\node[curveblue] at (axis cs:2.05,1.45) {$f$};
\node[targetgreen] at (axis cs:2.05,0.12) {$g$};
\node[black] at (axis cs:1.0,-0.32) {shared zero};
\end{axis}
\end{tikzpicture}
\caption{Two functions sharing a zero at $x=a$. Close to $a$ each is
indistinguishable from its tangent line (dashed). The quotient $f/g$ therefore behaves
like the quotient of the tangent slopes, $f'(a)/g'(a)$, and the troublesome common
factor $(x-a)$ cancels.}
\label{ch18:fig-rates}
\end{figure}

\begin{trapbox}
This tangent-line argument is the right \emph{intuition}, but it is not yet a proof,
and it quietly assumes more than L'Hopital's rule needs. It assumes $f$ and $g$ are
differentiable \emph{at} $a$ with $g'(a)\neq 0$, and it hides what happens when both
slopes are zero (giving $\tfrac{0}{0}$ all over again). The actual theorem, proved
below with the Mean Value Theorem, is sharper: it only needs the derivatives near $a$,
not the cancellation of tangent lines, and it lets you repeat the trick when the first
attempt stalls.
\end{trapbox}

\section{L'Hopital's rule, stated precisely}

\begin{theorem}[L'Hopital's rule, the $0/0$ case]\label{ch18:thm-lhopital}
Let $f$ and $g$ be differentiable on an open interval containing $a$ (except possibly
at $a$ itself), with $g'(x)\neq 0$ near $a$. Suppose
\[
\lim_{x\to a}f(x)=0\qquad\text{and}\qquad \lim_{x\to a}g(x)=0.
\]
If $\displaystyle\lim_{x\to a}\frac{f'(x)}{g'(x)}$ exists (as a finite number or as
$\pm\infty$), then
\[
\lim_{x\to a}\frac{f(x)}{g(x)}=\lim_{x\to a}\frac{f'(x)}{g'(x)}.
\]
\end{theorem}

Several features of this statement are doing real work, so read them deliberately.
The rule promises only one direction: \emph{if} the limit of the derivative ratio
exists, then the original limit exists and equals it. It says nothing if the
derivative ratio fails to settle down --- in that case the rule is simply silent, and
the original limit might still exist for another reason. Notice too that we never need
$f$ or $g$ to be defined at $a$; only their behavior \emph{near} $a$ matters, which is
exactly right for a limit. And the very same statement holds verbatim for one-sided
limits, for $x\to\pm\infty$, and --- with one change we will note --- for the
$\tfrac{\infty}{\infty}$ case.

\begin{whybox}
\textbf{Why the rule works: it is the Mean Value Theorem applied to two functions at
once.} The clean proof does not lean on tangent lines at all. It leans on
\emph{Cauchy's Mean Value Theorem}, the two-function version of the MVT you met in the
chapter on the Mean Value Theorem. Cauchy's theorem says that for suitable $f,g$ on an
interval $[a,x]$ there is a point $c$ strictly between $a$ and $x$ with
\[
\frac{f(x)-f(a)}{g(x)-g(a)}=\frac{f'(c)}{g'(c)}.
\]
That is the bridge. The left side is built from \emph{values}; the right side is built
from \emph{derivatives}. If $f(a)=g(a)=0$, the left side collapses to exactly
$\tfrac{f(x)}{g(x)}$ --- the thing we want --- and it has been rewritten as a ratio of
derivatives at some in-between point $c$. As $x\to a$, that point $c$ is squeezed
toward $a$ as well (it is trapped between $a$ and $x$), so the ratio
$\tfrac{f'(c)}{g'(c)}$ is forced to the limit of $\tfrac{f'}{g'}$. The values turn into
derivatives by the MVT, and the derivatives carry the answer.
\end{whybox}

\begin{focusbox}
\textbf{Proof (the $0/0$ case, finite $a$).}
First extend $f$ and $g$ to be continuous at $a$ by \emph{defining} $f(a)=g(a)=0$;
since $f,g\to 0$ as $x\to a$, this loses nothing and makes them continuous on a closed
interval with endpoint $a$. Fix $x$ near $a$, say $x>a$ (the case $x<a$ is identical).
On $[a,x]$ both functions are continuous and on $(a,x)$ both are differentiable, and
$g'\neq 0$ there. Cauchy's Mean Value Theorem then supplies a point $c\in(a,x)$ with
\[
\frac{f(x)-f(a)}{g(x)-g(a)}=\frac{f'(c)}{g'(c)}.
\]
(En route, $g(x)\neq g(a)$ because if they were equal the ordinary MVT would force
$g'$ to vanish somewhere in $(a,x)$, which we excluded; so the left side is legal.)
Using $f(a)=g(a)=0$, the left side is just $\tfrac{f(x)}{g(x)}$, hence
\[
\frac{f(x)}{g(x)}=\frac{f'(c)}{g'(c)},\qquad a<c<x.
\]
Now let $x\to a^+$. The point $c$ is pinned in $(a,x)$, so $c\to a^+$ too. By
hypothesis $\tfrac{f'(t)}{g'(t)}\to L$ as $t\to a$, so $\tfrac{f'(c)}{g'(c)}\to L$, and
therefore $\tfrac{f(x)}{g(x)}\to L$. The left-hand limit is the same argument on
$[x,a]$. Combining one-sided limits gives the theorem. $\qquad\blacksquare$
\end{focusbox}

\begin{remark}
The $\tfrac{\infty}{\infty}$ case is true with the identical conclusion, but its proof
is fussier --- you cannot simply set $f(a)=g(a)=0$ because the functions are blowing
up, not vanishing. The correct argument splits the interval and controls a tail term;
it is a standard but more technical $\eps$ estimate. The good news is that you may use
it with the same confidence and the same mechanics. We state it next and move on.
\end{remark}

\begin{theorem}[L'Hopital's rule, the $\infty/\infty$ case]\label{ch18:thm-lhopital-inf}
Under the same differentiability hypotheses, if $\;\abs{f(x)}\to\infty$ and
$\abs{g(x)}\to\infty$ as $x\to a$ (or as $x\to\pm\infty$), and if
$\lim \tfrac{f'(x)}{g'(x)}$ exists (finite or $\pm\infty$), then
$\lim \tfrac{f(x)}{g(x)}$ exists and equals it.
\end{theorem}

\begin{historybox}
The rule is named for the Marquis \textbf{Guillaume de l'H\^opital}, who published it
in 1696 in the first-ever calculus textbook, \emph{Analyse des Infiniment Petits}. But
he did not discover it. He was paying \textbf{Johann Bernoulli} a handsome retainer for
private lessons and for ``his'' mathematical discoveries, and the rule was Bernoulli's,
delivered by letter. L'Hopital was candid in his preface that he owed much to the
Bernoullis and to Leibniz, but the result went out under his name, and after his death
Bernoulli loudly --- and rightly --- claimed it. The surviving correspondence settles
the matter in Bernoulli's favor. So the most famous named rule in elementary calculus
is, by a quirk of patronage, named after the man who \emph{bought} it rather than the
man who \emph{found} it. The original 1696 statement was purely geometric, phrased in
the language of infinitesimals; the rigorous version resting on the Mean Value Theorem
had to wait for Cauchy in the 1820s, when limits were finally put on solid ground.
\end{historybox}

\section{Worked examples on the two ratio forms}

\begin{example}[The classic $\sin$ limit, again]\label{ch18:ex-sinx}
Consider $\displaystyle\lim_{x\to 0}\frac{\sin x}{x}$. Substituting gives
$\tfrac{0}{0}$. Differentiate top and bottom \emph{separately} (this is not the
quotient rule --- more on that trap shortly):
\[
\lim_{x\to 0}\frac{\sin x}{x}\overset{0/0}{=}\lim_{x\to 0}\frac{\cos x}{1}
=\frac{\cos 0}{1}=1.
\]
\end{example}

\begin{trapbox}
That last example is a delicious circular trap if you are being rigorous. We computed
$\ddx{\sin x}=\cos x$ in the first place by \emph{using}
$\lim_{x\to0}\tfrac{\sin x}{x}=1$. So ``proving'' that limit with L'Hopital is
logically circular --- it assumes what it sets out to show. L'Hopital gives the right
\emph{number}, but here it is not a legitimate proof. Use the geometric squeeze
argument for that one. The moral: a tool that gives the right answer can still be the
wrong argument.
\end{trapbox}

\begin{example}[A polynomial-flavored $0/0$]
$\displaystyle\lim_{x\to 1}\frac{x^3-1}{x^2-1}$ is $\tfrac{0}{0}$. By L'Hopital,
\[
\lim_{x\to 1}\frac{x^3-1}{x^2-1}\overset{0/0}{=}\lim_{x\to 1}\frac{3x^2}{2x}
=\frac{3}{2}.
\]
Sanity check by factoring: $\frac{x^3-1}{x^2-1}=\frac{(x-1)(x^2+x+1)}{(x-1)(x+1)}
=\frac{x^2+x+1}{x+1}\to\frac{3}{2}$. The factoring \emph{is} the cancellation of the
common $(x-1)$ that L'Hopital performs invisibly --- same mechanism, two faces.
\end{example}

\begin{example}[When once is not enough: applying the rule twice]\label{ch18:ex-twice}
Take $\displaystyle\lim_{x\to 0}\frac{1-\cos x}{x^2}$, which is $\tfrac{0}{0}$. One
application:
\[
\lim_{x\to 0}\frac{1-\cos x}{x^2}\overset{0/0}{=}\lim_{x\to 0}\frac{\sin x}{2x}.
\]
This new limit is \emph{still} $\tfrac{0}{0}$ --- the rule did not finish the job, but
it left us a simpler $\tfrac00$. So apply it again:
\[
\lim_{x\to 0}\frac{\sin x}{2x}\overset{0/0}{=}\lim_{x\to 0}\frac{\cos x}{2}
=\frac{1}{2}.
\]
You may iterate L'Hopital as many times as each successive limit stays a legal
$\tfrac00$ or $\tfrac{\infty}{\infty}$. The instant a step yields a determinate form,
stop and evaluate --- one more differentiation past that point will give a wrong
answer.
\end{example}

\begin{trapbox}
\textbf{Check the form before \emph{every} application.} The single most common
L'Hopital error is to differentiate when the form is no longer indeterminate. In the
example above, the moment we reached $\tfrac{\cos x}{2}$ the form became
$\tfrac{1}{2}$, a perfectly ordinary number; differentiating again would give
$\tfrac{-\sin x}{0}$, which is garbage. L'Hopital is a \emph{conditional} rule. Verify
$\tfrac00$ or $\tfrac{\infty}{\infty}$ at each step, every time, before you reach for a
derivative.
\end{trapbox}

\begin{labbox}
Compute $\displaystyle\lim_{x\to 0}\frac{e^x-1-x}{x^2}$ by hand. Confirm it is
$\tfrac00$. Apply L'Hopital once and check you again get $\tfrac00$. Apply it a second
time and read off the answer. (You should land on $\tfrac12$. Notice how this mirrors
the previous example --- both are detecting the \emph{quadratic} coefficient in a
Taylor expansion, a connection we make explicit at the end of the chapter.)
\end{labbox}

\begin{example}[An $\infty/\infty$ at infinity]\label{ch18:ex-lnx}
$\displaystyle\lim_{x\to\infty}\frac{\ln x}{x}$ is $\tfrac{\infty}{\infty}$. By
Theorem~\ref{ch18:thm-lhopital-inf},
\[
\lim_{x\to\infty}\frac{\ln x}{x}\overset{\infty/\infty}{=}
\lim_{x\to\infty}\frac{1/x}{1}=\lim_{x\to\infty}\frac{1}{x}=0.
\]
So $\ln x$, though it does march off to infinity, is utterly outrun by $x$. We will
build a whole hierarchy out of facts like this in the last section.
\end{example}

\section{Forcing the other forms into a fraction}

L'Hopital only speaks the language of fractions, $\tfrac00$ and
$\tfrac{\infty}{\infty}$. The remaining five indeterminate forms are not fractions on
their face. The art is to \emph{algebraically rewrite} each one as a fraction without
changing its value, and then let the rule loose. Each rewriting is a small, learnable
move.

\subsection{The product form $0\cdot\infty$}

A product $f\cdot g$ where $f\to 0$ and $g\to\infty$ is a tug-of-war: the $0$ tries to
crush the product to nothing, the $\infty$ tries to blow it up. Who wins is exactly the
kind of question L'Hopital answers --- once we make it a fraction. And there is a
fraction hiding in any product, because dividing is multiplying by a reciprocal:
\[
f\cdot g=\frac{f}{1/g}\quad\Big(\text{form }\tfrac00\Big)
\qquad\text{or}\qquad
f\cdot g=\frac{g}{1/f}\quad\Big(\text{form }\tfrac{\infty}{\infty}\Big).
\]
You get to choose which factor to flip. Flip the one whose reciprocal is friendlier to
differentiate.

\begin{example}[$x\ln x$ at the origin]\label{ch18:ex-xlnx}
$\displaystyle\lim_{x\to 0^+}x\ln x$ is $0\cdot(-\infty)$. Write it as a fraction by
flipping the $x$:
\[
x\ln x=\frac{\ln x}{1/x}\quad\Big(\tfrac{-\infty}{\infty}\Big)
\quad\Longrightarrow\quad
\lim_{x\to 0^+}\frac{\ln x}{1/x}\overset{\infty/\infty}{=}
\lim_{x\to 0^+}\frac{1/x}{-1/x^2}
=\lim_{x\to 0^+}\big(-x\big)=0.
\]
Had we flipped the other way, $\tfrac{x}{1/\ln x}$, the derivative of $1/\ln x$ is
$-\tfrac{1}{x(\ln x)^2}$ and the algebra is uglier. \emph{Choosing the right
reciprocal} is most of the skill in these problems.
\end{example}

\begin{walkbox}
Before reading on, decide on a walk: in $\lim_{x\to 0^+}x\ln x$, the function $x$ is
pulling the product to $0$ and $\ln x$ is pulling it to $-\infty$. The answer came out
$0$. So which pull is ``stronger,'' and what does that say about how fast $\ln x$ blows
up compared to how fast $x$ shrinks? Hold that thought until the growth-rates section.
\end{walkbox}

\subsection{The difference form $\infty-\infty$}

Here two quantities both race to infinity and we subtract them; the answer depends on
which is bigger and by how much. There is no universal recipe, but there are two
reliable moves: \emph{combine into a single fraction} (common denominator), or
\emph{factor out the dominant term}. Either one converts the subtraction into a ratio.

\begin{example}[A common-denominator rescue]\label{ch18:ex-cscdiff}
$\displaystyle\lim_{x\to 0}\left(\frac{1}{\sin x}-\frac{1}{x}\right)$ is
$\infty-\infty$ (each piece blows up). Put them over a common denominator:
\[
\frac{1}{\sin x}-\frac{1}{x}=\frac{x-\sin x}{x\sin x}.
\]
Now it is a clean $\tfrac00$. L'Hopital once:
\[
\lim_{x\to 0}\frac{x-\sin x}{x\sin x}
\overset{0/0}{=}\lim_{x\to 0}\frac{1-\cos x}{\sin x+x\cos x},
\]
still $\tfrac00$. Again:
\[
\overset{0/0}{=}\lim_{x\to 0}\frac{\sin x}{2\cos x-x\sin x}
=\frac{0}{2-0}=0.
\]
So the two infinities cancel \emph{exactly} in the limit; their difference settles to
$0$.
\end{example}

\subsection{The power forms $1^{\infty}$, $0^{0}$, $\infty^{0}$}

These three look different from each other but yield to one universal trick: a power
becomes a product the moment you take a logarithm. If $y=f(x)^{g(x)}$, then
\[
\ln y=g(x)\,\ln f(x),
\]
and each of the three power forms turns the right-hand side into a $0\cdot\infty$
product, which we already know how to fractionate. Find $\lim \ln y=L$, then
exponentiate: $\lim y=e^{L}$. The logarithm trades an unmanageable exponent for a
manageable product --- that is the whole idea.

\begin{whybox}
\textbf{Why the logarithm tames every power form.} The trouble in $1^\infty$ and its
cousins is that the \emph{exponent} is doing the limiting, and our tools (derivatives,
L'Hopital) are built for sums, products, and ratios --- not for moving exponents.
Taking $\ln$ is precisely the operation that \emph{pulls the exponent down} into a
coefficient: $\ln(f^{g})=g\ln f$. An impossible-looking power becomes an ordinary
product; the product becomes a fraction; the fraction yields to L'Hopital. Because
$\exp$ is continuous, $\lim \ln y=L$ guarantees $\lim y=e^{L}$, so nothing is lost in
the round trip $y\to\ln y\to L\to e^L$. We are not avoiding the exponent --- we are
\emph{lowering} it to where our machinery can reach.
\end{whybox}

\begin{example}[The most famous $1^{\infty}$: compound interest revisited]\label{ch18:ex-e}
$\displaystyle\lim_{x\to\infty}\left(1+\frac{1}{x}\right)^{x}$ is $1^{\infty}$ --- the
base creeps to $1$, the exponent runs to $\infty$, and they fight. Let
$y=\left(1+\tfrac1x\right)^{x}$ and take logs:
\[
\ln y=x\ln\!\left(1+\frac1x\right).
\]
This is $\infty\cdot 0$. Make it a fraction by flipping the $x$:
\[
\ln y=\frac{\ln\!\left(1+\tfrac1x\right)}{1/x}\quad\Big(\tfrac00\Big).
\]
Differentiate top and bottom. The top differentiates by the chain rule:
$\ddx{\ln(1+1/x)}=\frac{1}{1+1/x}\cdot\left(-\tfrac{1}{x^2}\right)$, and the bottom
gives $-\tfrac{1}{x^2}$:
\[
\lim_{x\to\infty}\frac{\dfrac{-1/x^2}{1+1/x}}{-1/x^2}
=\lim_{x\to\infty}\frac{1}{1+1/x}=1.
\]
So $\lim \ln y=1$, hence $\lim y=e^{1}=e$. This recovers, by L'Hopital, the very
limit-definition of $e$ from the early chapters --- the compound-interest number falls
right out. Figure~\ref{ch18:fig-e} shows the approach.
\end{example}

\begin{figure}[H]
\centering
\begin{tikzpicture}
\begin{axis}[
  width=11cm, height=6cm,
  axis lines=middle,
  xlabel={$x$}, ylabel={$y$},
  xmin=0, xmax=30, ymin=1.8, ymax=2.9,
  samples=120, domain=0.5:30,
  legend pos=south east,
  clip=false,
]
\addplot[curveblue, very thick] {exp(x*ln(1+1/x))};
\addlegendentry{$y=\left(1+\tfrac1x\right)^x$}
\addplot[tangentred, dashed, domain=0:30] {2.718281828};
\addlegendentry{$y=e$}
\end{axis}
\end{tikzpicture}
\caption{The compound-interest sequence $\left(1+\tfrac1x\right)^x$ climbing toward
$e\approx 2.71828$ as $x\to\infty$. The form is $1^\infty$ --- the base $1+\tfrac1x$
sinks to $1$ while the exponent $x$ grows --- and the two effects balance at exactly
$e$. The dashed line is the limit the curve approaches but never reaches.}
\label{ch18:fig-e}
\end{figure}

\begin{example}[An $\infty^{0}$]\label{ch18:ex-inf0}
$\displaystyle\lim_{x\to\infty}x^{1/x}$ is $\infty^{0}$. Logs:
$\ln y=\tfrac{\ln x}{x}$, which we already found in Example~\ref{ch18:ex-lnx} to
be $\tfrac{\infty}{\infty}\to 0$. Therefore $\lim y=e^{0}=1$. A number raised to a
gently vanishing power slides back down to $1$.
\end{example}

\begin{example}[A $0^{0}$]\label{ch18:ex-00}
$\displaystyle\lim_{x\to 0^+}x^{x}$ is $0^{0}$, perhaps the most counterintuitive of
the lot. Logs: $\ln y=x\ln x$, which is the $0\cdot(-\infty)$ from
Example~\ref{ch18:ex-xlnx}, equal to $0$. Hence $\lim_{x\to0^+}x^x=e^{0}=1$. So $x^x$,
despite the base sinking to $0$, climbs back up to $1$. The figure makes this vivid.
\end{example}

\begin{figure}[H]
\centering
\begin{tikzpicture}
\begin{axis}[
  width=10.5cm, height=6cm,
  axis lines=middle,
  xlabel={$x$}, ylabel={$y$},
  xmin=0, xmax=1.5, ymin=0, ymax=1.5,
  samples=120, domain=0.02:1.5,
  legend pos=south east,
  clip=false,
]
\addplot[curveblue, very thick] {exp(x*ln(x))};
\addlegendentry{$y=x^x$}
\addplot[only marks, mark=*, mark size=1.6pt, tangentred] coordinates {(0.02,0.925)};
\node[tangentred] at (axis cs:0.55,1.05) {heads to $1$ as $x\to 0^+$};
\addplot[gray, dashed, domain=0:1.5] {1};
\end{axis}
\end{tikzpicture}
\caption{The graph of $y=x^x$ for $x>0$. Although the base $x$ tends to $0$, the value
$x^x$ does not collapse to $0$; it dips to a minimum near $x\approx 0.368$ and
approaches $1$ as $x\to 0^+$. The form $0^0$ is genuinely indeterminate, and here it
resolves to $1$.}
\label{ch18:fig-xx}
\end{figure}

\begin{trapbox}
$0^0$, $1^\infty$, and $\infty^0$ are \emph{forms}, not values. People sometimes
``know'' that $0^0=1$ from algebra (where the convention is handy) and conclude that
$\lim x^x$ ``must'' be $1$ for that reason. The reasoning is wrong even though the
answer is right: change the problem to $\lim_{x\to 0^+}x^{1/\ln x}$, also $0^0$, and
the answer is $e$, not $1$. The form does not decide; only the
logarithm-and-L'Hopital computation does.
\end{trapbox}

\section{How the rule gets misused}

The rule is convenient enough that students reach for it reflexively, and that reflex
causes most of the wrong answers. Here are the failure modes, each worth memorizing as
a thing \emph{not} to do.

\begin{trapbox}
\textbf{Trap 1: Using it on a determinate form.} If the form is not $\tfrac00$ or
$\tfrac{\infty}{\infty}$, the rule simply does not apply, and using it gives nonsense.
For instance $\displaystyle\lim_{x\to 0}\frac{\cos x}{x+1}$ is $\tfrac{1}{1}=1$ by
direct substitution. ``Applying L'Hopital'' would give
$\lim \tfrac{-\sin x}{1}=0$ --- flatly wrong. \emph{Always check the form first.}
\end{trapbox}

\begin{trapbox}
\textbf{Trap 2: Confusing it with the quotient rule.} L'Hopital says differentiate the
numerator and the denominator \emph{separately} and form a new quotient
$\tfrac{f'}{g'}$. It does \emph{not} say differentiate the quotient $\tfrac{f}{g}$ with
the quotient rule $\tfrac{f'g-fg'}{g^2}$. These are completely different expressions.
Writing the quotient rule here is one of the most frequent and most fatal slips.
\end{trapbox}

\begin{trapbox}
\textbf{Trap 3: Continuing after the form resolves.} As stressed in
Example~\ref{ch18:ex-twice}, the instant a step is no longer indeterminate you must
stop. One extra differentiation past a determinate form gives a wrong number with
total confidence.
\end{trapbox}

\begin{trapbox}
\textbf{Trap 4: Treating a non-existent derivative limit as ``no answer.''} The rule is
one-directional. If $\lim\tfrac{f'}{g'}$ fails to exist, you may conclude
\emph{nothing} about $\lim\tfrac{f}{g}$; it might still exist. Take
$\displaystyle\lim_{x\to\infty}\frac{x+\sin x}{x}$. This is $\tfrac{\infty}{\infty}$,
and L'Hopital gives $\lim \tfrac{1+\cos x}{1}=\lim(1+\cos x)$, which \emph{oscillates}
and has no limit. A beginner concludes ``the limit does not exist.'' But the original
limit plainly \emph{does} exist: $\frac{x+\sin x}{x}=1+\frac{\sin x}{x}\to 1+0=1$.
L'Hopital was silent here, not authoritative.
\end{trapbox}

\begin{trapbox}
\textbf{Trap 5: An endless loop.} $\displaystyle\lim_{x\to\infty}\frac{e^x}{e^x+1}$ is
$\tfrac{\infty}{\infty}$; L'Hopital gives $\tfrac{e^x}{e^x}=1$ in one clean step, fine.
But $\displaystyle\lim_{x\to\infty}\frac{\sqrt{x^2+1}}{x}$, attacked by L'Hopital,
produces $\frac{x/\sqrt{x^2+1}}{1}=\frac{x}{\sqrt{x^2+1}}$, whose own L'Hopital sends
you straight back to $\frac{\sqrt{x^2+1}}{x}$. You are chasing your tail. Step back and
use algebra: $\frac{\sqrt{x^2+1}}{x}=\sqrt{1+\tfrac{1}{x^2}}\to 1$. When L'Hopital
cycles, abandon it.
\end{trapbox}

\begin{recoverybox}
A healthy habit before reaching for L'Hopital: ask whether a quick algebraic move ---
factoring, a common denominator, dividing by the dominant power, a known limit ---
settles the problem first. L'Hopital is a hammer; not every limit is a nail. The
examples in Traps 4 and 5 are solved faster, and more safely, by elementary algebra.
\end{recoverybox}

\section{Growth rates, settled once and for all}

We have brushed up against a recurring question since the chapter on limits at
infinity: among the functions that all run off to infinity --- powers of $x$,
exponentials, logarithms --- \emph{which outruns which}? L'Hopital lets us answer this
cleanly, because ``$f$ grows faster than $g$'' is precisely the statement
\[
\lim_{x\to\infty}\frac{g(x)}{f(x)}=0,
\]
an $\tfrac{\infty}{\infty}$ form tailor-made for the rule.

\begin{theorem}[The growth hierarchy]\label{ch18:thm-growth}
As $x\to\infty$, for any fixed exponents $a>0$, $p>0$ and any base $b>1$,
\[
(\ln x)^{a}\;\ll\; x^{p}\;\ll\; b^{x}\;\ll\; x!\;\ll\; x^{x},
\]
where $u\ll v$ means $\lim_{x\to\infty}\tfrac{u}{v}=0$ (read: ``$v$ dominates $u$'').
In words: every power of a logarithm loses to every positive power of $x$; every
positive power of $x$ loses to every exponential with base above $1$; the exponential
loses to the factorial; and the factorial loses to $x^x$.
\end{theorem}

\begin{focusbox}
\textbf{Power beats logarithm.} We showed $\tfrac{\ln x}{x}\to 0$. More generally
$\tfrac{\ln x}{x^p}$ for any $p>0$ is $\tfrac{\infty}{\infty}$, and L'Hopital gives
$\tfrac{1/x}{p\,x^{p-1}}=\tfrac{1}{p\,x^{p}}\to 0$. So $\ln x\ll x^p$ for every $p>0$,
no matter how tiny --- even $x^{0.001}$ eventually buries $\ln x$. Raising the log to a
power $a$ only delays the inevitable: $a$ applications of L'Hopital grind
$(\ln x)^a/x^p$ down to $0$ as well.
\end{focusbox}

\begin{focusbox}
\textbf{Exponential beats power.} Consider $\tfrac{x^n}{e^x}$ for a positive integer
$n$ --- an $\tfrac{\infty}{\infty}$. Each L'Hopital step lowers the power on top by one
while the $e^x$ on the bottom regenerates unchanged:
\[
\frac{x^n}{e^x}\to\frac{n x^{n-1}}{e^x}\to\frac{n(n-1)x^{n-2}}{e^x}\to\cdots
\to\frac{n!}{e^x}\to 0.
\]
After $n$ steps the numerator is the constant $n!$ and the denominator still races to
infinity, so the ratio dies. This is the precise sense in which exponential growth
\emph{always} overwhelms polynomial growth --- the engine of why compound processes,
runaway populations, and $2^n$-time algorithms eventually dwarf any polynomial.
\end{focusbox}

\begin{figure}[H]
\centering
\begin{tikzpicture}
\begin{axis}[
  width=11cm, height=6.5cm,
  axis lines=middle,
  xlabel={$x$}, ylabel={$y$},
  xmin=0, xmax=6, ymin=0, ymax=60,
  samples=120, domain=0.05:6,
  legend pos=north west,
  clip=false,
]
\addplot[targetgreen, very thick] {ln(x)};
\addlegendentry{$\ln x$}
\addplot[curveblue, very thick] {x^2};
\addlegendentry{$x^2$}
\addplot[tangentred, very thick, domain=0.05:4.09] {exp(x)};
\addlegendentry{$e^x$}
\end{axis}
\end{tikzpicture}
\caption{The growth hierarchy in one window. $\ln x$ barely lifts off; $x^2$ climbs
steadily; $e^x$ is already running away by $x\approx 4$. Pushed to infinity, the order
is permanent: $\ln x\ll x^2\ll e^x$. L'Hopital's rule is what proves the order never
reverses.}
\label{ch18:fig-growth}
\end{figure}

\begin{applicationbox}
\textbf{Why your sorting algorithm scales but your brute-force search does not.} In
computer science, an algorithm's running time as a function of input size $n$ is
exactly a growth-rate question. A method costing $n\log n$ steps (good sorting) versus
one costing $n^2$ (naive sorting) versus one costing $2^n$ (brute-force subset search)
are separated by precisely the hierarchy in Theorem~\ref{ch18:thm-growth}:
$n\log n\ll n^2\ll 2^n$. The L'Hopital computations above are the rigorous reason a
$2^n$ algorithm becomes hopeless near $n=60$ while an $n\log n$ algorithm sails through
$n$ in the billions. Growth-rate order is not a heuristic; it is a theorem, and this is
its proof.
\end{applicationbox}

\begin{connectionbox}
The factorial and $x^x$ ends of the hierarchy connect to \emph{Stirling's
approximation}, $n!\sim \sqrt{2\pi n}\,(n/e)^n$, which you will meet in probability and
combinatorics. Stirling shows $n!$ sits, roughly, like $(n/e)^n$ --- slightly below
$n^n$, comfortably above $b^n$ --- which is exactly the placement
$b^x\ll x!\ll x^x$ in our chain. The crude L'Hopital ordering and the refined Stirling
estimate tell the same story at two levels of resolution.
\end{connectionbox}

\section{The deeper picture: L'Hopital is Taylor's theorem in miniature}

The tangent-line intuition we opened with was not a coincidence. When you push it one
order further, L'Hopital's rule is revealed as a low-resolution shadow of \emph{Taylor
series}, which you will meet in full later. The idea: near a zero, a function looks not
just like its tangent line but like its whole power series, and the limit of a ratio is
decided by the \emph{first surviving terms} of the top and bottom.

\begin{happeningbox}
\textbf{Why repeated L'Hopital keeps working.} Suppose near $x=0$,
\[
f(x)=\frac{f''(0)}{2}x^2+\cdots,\qquad g(x)=g'(0)\,x+\cdots
\]
Then $\tfrac{f}{g}\sim \tfrac{(f''(0)/2)x^2}{g'(0)x}=\tfrac{f''(0)}{2g'(0)}x\to 0$.
The function whose series \emph{starts at a higher power} vanishes faster, so it loses.
Each L'Hopital step differentiates away one power of $x$ from both series; you keep
going until the surviving lowest-order terms no longer both vanish, and at that moment
the ratio of their coefficients is the limit. That is why
$\tfrac{1-\cos x}{x^2}\to\tfrac12$: the series
$1-\cos x=\tfrac{x^2}{2}-\tfrac{x^4}{24}+\cdots$ starts at $\tfrac{x^2}{2}$, exactly
matching the $x^2$ below, and the ratio of leading coefficients is
$\tfrac{1/2}{1}=\tfrac12$ --- the same $\tfrac12$ two L'Hopital steps produced. The
rule is reading off the first nonzero coefficients in disguise.
\end{happeningbox}

\begin{aestheticbox}
There is something satisfying in seeing one idea reappear in three costumes. The
tangent-line picture, the Cauchy MVT proof, and the Taylor-series view are not three
facts about L'Hopital's rule --- they are one fact at three depths. The tangent line is
the first-order truth; the MVT makes it rigorous without assuming differentiability at
the point; the Taylor series explains why iterating the rule must terminate and what
number it terminates at. When a theorem looks the same from every angle you approach
it, you have understood it.
\end{aestheticbox}

\section{Exercises}

\begin{exercise}[Warm-up ratios]
Evaluate, checking the form before each L'Hopital step:
(a) $\displaystyle\lim_{x\to 0}\frac{\tan x - x}{x^3}$;\quad
(b) $\displaystyle\lim_{x\to 1}\frac{\ln x}{x-1}$;\quad
(c) $\displaystyle\lim_{x\to\infty}\frac{x^2}{e^{x}}$;\quad
(d) $\displaystyle\lim_{x\to 0}\frac{e^{x}-1}{\sin x}$.
\end{exercise}

\begin{exercise}[Products and differences]
Convert to a fraction first, then evaluate:
(a) $\displaystyle\lim_{x\to 0^+}x^2\ln x$;\quad
(b) $\displaystyle\lim_{x\to (\pi/2)^-}(\sec x - \tan x)$;\quad
(c) $\displaystyle\lim_{x\to 0}\left(\frac{1}{x}-\frac{1}{e^x-1}\right)$.
\end{exercise}

\begin{exercise}[Power forms]
Use the logarithm trick:
(a) $\displaystyle\lim_{x\to\infty}\left(1+\frac{3}{x}\right)^{x}$;\quad
(b) $\displaystyle\lim_{x\to 0^+}(\sin x)^{x}$;\quad
(c) $\displaystyle\lim_{x\to\infty}\left(\frac{x+1}{x-1}\right)^{x}$.
\end{exercise}

\begin{exercise}[Spot the misuse]
For each, decide whether L'Hopital is legitimately applicable, and if not, find the
limit by other means.
(a) $\displaystyle\lim_{x\to 0}\frac{x+2}{x^2+1}$;\quad
(b) $\displaystyle\lim_{x\to\infty}\frac{2x-\cos x}{x+1}$;\quad
(c) $\displaystyle\lim_{x\to\infty}\frac{\sqrt{4x^2+1}}{x+3}$.
In (b) and (c), explain in one sentence why blindly differentiating fails or stalls.
\end{exercise}

\begin{exercise}[Three steps, then stop]
Evaluate $\displaystyle\lim_{x\to 0}\frac{x-\sin x}{x^3}$. State the form at each stage
and identify the exact moment you must stop applying the rule. Confirm your answer
against the series $\sin x = x-\tfrac{x^3}{6}+\cdots$.
\end{exercise}

\begin{exercise}[Growth-rate ordering]
Order the following by eventual size as $x\to\infty$, smallest first, and justify each
pairwise comparison with a single L'Hopital limit:
$\;x^{10},\;\; (\ln x)^{100},\;\; 1.0001^{\,x},\;\; e^{\sqrt{x}}.$
(The last one is subtler than it looks --- compare its logarithm.)
\end{exercise}

\begin{exercise}[A famous compound limit]
Show that $\displaystyle\lim_{x\to\infty}\left(1-\frac{r}{x}\right)^{x}=e^{-r}$ for any
constant $r$, by taking logs and applying L'Hopital. Interpret the result for
continuously compounded \emph{decay} (e.g.\ radioactive material or a depreciating
asset losing a fixed fraction per ever-shorter period).
\end{exercise}

\begin{exercise}[Conceptual: why the form matters]
Construct two functions $f,g$ with $f,g\to 0$ as $x\to 0$ for which
$\lim\tfrac{f}{g}=7$, and another pair for which the same form $\tfrac00$ gives
$\lim\tfrac{f}{g}=\infty$. Use your examples to explain, in your own words, the precise
sense in which $\tfrac00$ ``contains no information.''
\end{exercise}

\begin{exercise}[Proof: the simple case from the definition]
Suppose $f$ and $g$ are differentiable \emph{at} $a$ with $f(a)=g(a)=0$ and
$g'(a)\neq 0$. Without invoking the full theorem, prove directly from the definition of
the derivative that
\[
\lim_{x\to a}\frac{f(x)}{g(x)}=\frac{f'(a)}{g'(a)}.
\]
\emph{Hint:} divide numerator and denominator by $(x-a)$ and recognize two difference
quotients. Explain why this short argument needs the stronger hypothesis
``differentiable \emph{at} $a$,'' whereas Theorem~\ref{ch18:thm-lhopital} gets by with
less.
\end{exercise}

\begin{exercise}[Proof: L'Hopital can fail without the form]
Give an explicit pair $f,g$ with $\lim_{x\to a}\tfrac{f(x)}{g(x)}$ existing but
$\lim_{x\to a}\tfrac{f'(x)}{g'(x)}$ \emph{not} existing, where the form is genuinely
$\tfrac{\infty}{\infty}$. (Trap~4 contains the seed of such an example.) Explain how
this is consistent with the theorem --- that is, why no contradiction arises.
\end{exercise}
